\documentclass[11pt]{article}
\usepackage[T1]{fontenc}
\usepackage[utf8]{inputenc}
\usepackage{lmodern,amsmath,amssymb,amsthm}
\usepackage[margin=28mm]{geometry}
\usepackage[colorlinks=true,linkcolor=blue,urlcolor=blue]{hyperref}
\newtheorem{theorem}{Theorem}
\newtheorem{lemma}[theorem]{Lemma}
\theoremstyle{definition}
\newtheorem{definition}[theorem]{Definition}
\newtheorem{question}[theorem]{Open question}
\title{Classical Complexity Classes}
\author{\'{E}douard Bonnet \and Codex 5.6 and 6}
\date{}
\begin{document}
\maketitle
\begin{abstract}
We define classical complexity classes for languages of finite binary strings,
using explicit time bounds, work-space bounds, and certificate encodings.
The definition of $\mathrm P$ and its proved finite-stack and single-tape
characterizations are included locally. We prove
$\mathrm L\subseteq\mathrm{NL}\subseteq\mathrm P\subseteq\mathrm{NP}\subseteq\mathrm{PSPACE}=\mathrm{NPSPACE}\subseteq\mathrm{EXPTIME}$,
as well as the stated encoding properties and closure of $\mathrm P$ under
complement.
\end{abstract}

\section{Languages, machines, and complexity classes}
All languages are subsets of the finite binary strings. The following
definition text is reproduced from the concept layer. The annotations link
it to the precise machine data, encodings, quantifiers, and bounds in Lean.

\begin{definition}[The complexity class P]\leavevmode
% lax begin Lax434930.PolynomialTime
A language of finite binary strings belongs to $\mathrm{P}$ if a deterministic
Turing machine decides membership in that language in polynomial time.
Precisely, there are a single machine and a polynomial $p \in \mathbb{N}[X]$
such that, on every input $w$, the machine halts within $p(|w|)$ steps and
returns the bit $1$ if $w$ belongs to the language and $0$ otherwise.

The definition uses mathlib's deterministic stack machines, the identity
encoding of binary strings, and a singleton Boolean output. Time counts
transitions of fixed finite instruction blocks. We also prove equivalence
with elementary single-tape machines.
% lax end Lax434930.PolynomialTime
\end{definition}

\begin{definition}[Finite-stack and elementary single-tape definitions of P]\leavevmode
% lax begin Lax434930.MachineModels
The finite-stack class requires every stack alphabet to be finite. The
single-tape class uses a deterministic machine whose transitions each move
the head one square or write one symbol. Its alphabet and control states
are finite. The input is
written in its original order, starting at the head, with blank tape elsewhere.
The two input symbols are distinct and different from blank. A halted control
state determines the Boolean answer; work tape need not be erased.
% lax end Lax434930.MachineModels
\end{definition}

\begin{definition}[Binary encoding of an input and a certificate]\leavevmode
% lax begin Lax434930.Certificates
To encode a pair $(x,y)$ of binary strings, replace each bit $b$ of $x$ by
$0b$, then append a single $1$ followed by $y$. This encoding has length
$2|x|+|y|+1$ and has a unique decoding. In particular, a polynomial bound in
the encoded length is a polynomial bound in the combined input and
certificate lengths.
% lax end Lax434930.Certificates
\end{definition}

\begin{definition}[The complexity class NP]\leavevmode
% lax begin Lax434930.NondeterministicPolynomialTime
A binary language $A$ belongs to $\mathrm{NP}$ if there are a verifier
language $V\in\mathrm{P}$ and a polynomial $p\in\mathbb{N}[X]$ such that
$x\in A$ precisely when some binary certificate $y$ with $|y|\le p(|x|)$
satisfies $\langle x,y\rangle\in V$.
The verifier is a single deterministic polynomial-time decider on the
explicit pair encoding, and its running time is polynomial in the combined
input and certificate lengths. The certificate bound depends only on the
original input length. This is the standard certificate definition of NP.
% lax end Lax434930.NondeterministicPolynomialTime
\end{definition}

\begin{definition}[Finite Turing machines with a read-only input tape]\leavevmode
% lax begin Lax434930.SpaceMachines
A machine has finitely many control states, a finite work alphabet with a
blank symbol, a read-only binary input tape between distinct endmarkers,
and one initially blank, semi-infinite work tape. Both heads start at
position zero, which is the input's left endmarker and the work tape's
leftmost cell. One transition reads the two scanned symbols, writes one
work symbol, changes state, and moves each head by at most one cell.
An outward move at a tape boundary leaves that head in place.

The finite transition table gives a finite set of possible actions for
each state and pair of scanned symbols. A deterministic machine has at
most one action in every such set. A configuration with no successor is
terminal, and its control state supplies an accept/reject bit. Acceptance
means that some computation branch reaches an accepting terminal
configuration. A decider has a finite bound on the lengths of all branches
for each input, and accepts exactly the strings in its language. This
bound need not be computable or satisfy any specified time bound.

Only the work tape is charged as space. Bounding its head below $s$ at
every reachable configuration bounds every visited cell to the prefix
$0,\ldots,s-1$, including blank cells and cells later erased. The input
head is confined to $|x|+2$ positions and cannot act as an unbounded free
counter. The transition table sees neither head position nor the input
length; it sees only the control state and scanned symbols.
% lax end Lax434930.SpaceMachines
\end{definition}

\begin{definition}[Deterministic and nondeterministic space bounds]\leavevmode
% lax begin Lax434930.SpaceBounds
For a function $s:\mathbb{N}\to\mathbb{N}$, the classes DSPACE and NSPACE
below use the exact bound $s(n)$ on the number of work cells visited.
There is one machine for all inputs, it halts on every branch, and every
reachable configuration on an input of length $n$ respects the bound.
DSPACE additionally requires deterministic transitions. Constant factors
are quantified explicitly when the classical space classes are defined.
We use $\lfloor\log_2(n+2)\rfloor$ as a positive logarithmic bound, so
empty inputs are included without a special case.
% lax end Lax434930.SpaceBounds
\end{definition}

\begin{definition}[The complexity class L]\leavevmode
% lax begin Lax434930.LogarithmicSpace
A binary language belongs to $\mathrm{L}$ if a deterministic Turing machine
decides membership using $O(\log n)$ work space and a separate read-only
input tape. Precisely, some positive constant $c$ bounds work space by
$c\lfloor\log_2(n+2)\rfloor$ on every input of length $n$.
% lax end Lax434930.LogarithmicSpace
\end{definition}

\begin{definition}[The complexity class NL]\leavevmode
% lax begin Lax434930.NondeterministicLogarithmicSpace
A binary language belongs to $\mathrm{NL}$ if a nondeterministic Turing
machine decides membership using $O(\log n)$ work space and a separate
read-only input tape. Every branch halts and respects the space bound;
membership means that at least one branch accepts. The bound is
$c\lfloor\log_2(n+2)\rfloor$ for one positive constant $c$.
% lax end Lax434930.NondeterministicLogarithmicSpace
\end{definition}

\begin{definition}[The complexity class PSPACE]\leavevmode
% lax begin Lax434930.PolynomialSpace
A binary language belongs to $\mathrm{PSPACE}$ if one deterministic Turing
machine decides membership using at most $p(n)$ work cells on inputs of
length $n$, for some polynomial $p\in\mathbb{N}[X]$. The input tape is
read-only and excluded from work space, and the machine always halts.
% lax end Lax434930.PolynomialSpace
\end{definition}

\begin{definition}[The complexity class NPSPACE]\leavevmode
% lax begin Lax434930.NondeterministicPolynomialSpace
A binary language belongs to $\mathrm{NPSPACE}$ if one nondeterministic
Turing machine decides membership using at most $p(n)$ work cells on
every branch on inputs of length $n$, for some polynomial
$p\in\mathbb{N}[X]$. All branches halt; at least one accepts precisely
when the input belongs to the language. The input tape is read-only and
excluded from work space.
% lax end Lax434930.NondeterministicPolynomialSpace
\end{definition}

\begin{definition}[The complexity classes coNL and coNP]\leavevmode
% lax begin Lax434930.ComplementClasses
For a class $\mathcal C$ of binary languages, $\mathrm{co}\mathcal C$
consists of languages whose complements belong to $\mathcal C$.
Complements are taken among all finite binary strings. In particular,
$A\in\mathrm{coNL}$ means $\overline A\in\mathrm{NL}$, and
$A\in\mathrm{coNP}$ means $\overline A\in\mathrm{NP}$.
This operation complements each language; it does not take the set-theoretic
complement of the class of languages.
% lax end Lax434930.ComplementClasses
\end{definition}

\begin{definition}[The complexity class EXPTIME]\leavevmode
% lax begin Lax434930.ExponentialTime
A binary language belongs to $\mathrm{EXPTIME}$ if a deterministic Turing
machine decides membership within $2^{p(n)}$ transitions on every input
of length $n$, for some polynomial $p\in\mathbb{N}[X]$.
The machine is the finite elementary single-tape machine defined here:
input bits are distinct from blank, input appears in its original order,
and each transition performs one move or one write. A terminal state's
Boolean label gives the answer. The polynomial in the exponent may have
any fixed degree; this is the usual EXPTIME, also called EXP.
% lax end Lax434930.ExponentialTime
\end{definition}

\section{Encoding properties}

\begin{lemma}\leavevmode
% lax begin Lax434930.Certificates
Decoding $\langle x,y\rangle$ returns exactly $(x,y)$.
% lax end Lax434930.Certificates
\end{lemma}

\begin{proof}\leavevmode
% lax begin Lax434930Proofs.Certificates.unpair_pair
Induct on $x$. The empty word contributes only the delimiter; each additional bit contributes one two-bit block, which the decoder removes before applying the induction hypothesis.
\qedhere
% lax end Lax434930Proofs.Certificates.unpair_pair
\end{proof}

\begin{lemma}\leavevmode
% lax begin Lax434930.Certificates
The function $(x,y)\mapsto\langle x,y\rangle$ is injective.
% lax end Lax434930.Certificates
\end{lemma}

\begin{proof}\leavevmode
% lax begin Lax434930Proofs.Certificates.pair_injective
Apply the decoder to an equality of encodings and use the preceding decoding identity.
\qedhere
% lax end Lax434930Proofs.Certificates.pair_injective
\end{proof}

\begin{lemma}\leavevmode
% lax begin Lax434930.Certificates
The encoded pair has length $|\langle x,y\rangle|=2|x|+|y|+1$.
% lax end Lax434930.Certificates
\end{lemma}

\begin{proof}\leavevmode
% lax begin Lax434930Proofs.Certificates.pair_length
Induct on $x$. The delimiter contributes one bit, each bit of $x$ contributes two, and $y$ is appended unchanged.
\qedhere
% lax end Lax434930Proofs.Certificates.pair_length
\end{proof}

\section{Polynomial-time machine models and complement}

\begin{lemma}[Finite stack alphabets suffice]\leavevmode
% lax begin Lax434930.FiniteStackEquivalence
Requiring every work-stack alphabet to be finite leaves $\mathrm{P}$ unchanged.
% lax end Lax434930.FiniteStackEquivalence
\end{lemma}

\begin{proof}\leavevmode
% lax begin Lax434930Proofs.ModelEquivalence.finiteStackP_eq_P
The input and output alphabets are finite. The machine has finite control, and the push instructions have finite ranges. Restrict each stack alphabet to these symbols. Every reachable transition and its time bound are preserved. Forgetting the finite-alphabet restriction gives the reverse inclusion.
\qedhere
% lax end Lax434930Proofs.ModelEquivalence.finiteStackP_eq_P
\end{proof}

\begin{lemma}[Single-tape characterization of P]\leavevmode
% lax begin Lax434930.ModelEquivalence
The elementary single-tape and stack-machine definitions of $\mathrm{P}$
coincide. Both simulations include polynomial bounds for input conversion,
execution, and final output conversion.
% lax end Lax434930.ModelEquivalence
\end{lemma}

\begin{proof}\leavevmode
% lax begin Lax434930Proofs.ModelEquivalence.singleTapeP_eq_P
After restricting stack alphabets, compile the finitely many stacks to tape tracks with polynomial overhead. In the reverse direction, two stacks represent the portions of a tape on either side of its head, with linear overhead per simulated transition. The formal simulations include bounds for input conversion and final output conversion.
\qedhere
% lax end Lax434930Proofs.ModelEquivalence.singleTapeP_eq_P
\end{proof}

\begin{lemma}[P is closed under complement]\leavevmode
% lax begin Lax434930.ComplementClosure
If a language of binary strings belongs to $\mathrm{P}$, then its complement
also belongs to $\mathrm{P}$. The complement is taken in the set of all finite
binary strings.
% lax end Lax434930.ComplementClosure
\end{lemma}

\begin{proof}\leavevmode
% lax begin Lax434930Proofs.closed_under_complement
Compose the output-alphabet equivalence with Boolean negation. The unchanged execution then returns the opposite Boolean answer with the same polynomial bound. Negating the correctness equivalence identifies the complementary language.
\qedhere
% lax end Lax434930Proofs.closed_under_complement
\end{proof}

\begin{lemma}[Single-tape P is closed under complement]\leavevmode
% lax begin Lax434930.SingleTapeComplement
The elementary single-tape class $\mathrm{P}$ is closed under complement.
% lax end Lax434930.SingleTapeComplement
\end{lemma}

\begin{proof}\leavevmode
% lax begin Lax434930Proofs.ModelEquivalence.singleTape_closed_under_complement
Transport complement closure across the proved equality of the single-tape and stack-machine classes.
\qedhere
% lax end Lax434930Proofs.ModelEquivalence.singleTape_closed_under_complement
\end{proof}

\section{Inclusions}
\begin{theorem}[The inclusion chain]\leavevmode
% lax begin Lax434930.BasicProperties
The classes satisfy
\[
\mathrm{L}\subseteq\mathrm{NL}\subseteq\mathrm{P}\subseteq\mathrm{NP}
\subseteq\mathrm{PSPACE}=\mathrm{NPSPACE}\subseteq\mathrm{EXPTIME}.
\]
% lax end Lax434930.BasicProperties
\end{theorem}
The following six statements prove each inclusion and the equality.

\begin{lemma}\leavevmode
% lax begin Lax434930.BasicProperties
$\mathrm L\subseteq\mathrm{NL}$.
% lax end Lax434930.BasicProperties
\end{lemma}

\begin{proof}\leavevmode
% lax begin Lax434930Proofs.BasicProperties.L_subset_NL
Use the same machine and the same logarithmic space bound. Removing the requirement that each local observation has at most one action allows nondeterminism without changing any run.
\qedhere
% lax end Lax434930Proofs.BasicProperties.L_subset_NL
\end{proof}

\begin{lemma}\leavevmode
% lax begin Lax434930.BasicProperties
$\mathrm{NL}\subseteq\mathrm P$.
% lax end Lax434930.BasicProperties
\end{lemma}

\begin{proof}\leavevmode
% lax begin Lax434930Proofs.NL_subset_P
Fix a nondeterministic decider using at most $c\log_2(n+2)$ work cells.
Encode its state, input head, work head, and work-tape contents by a single
integer. Encoding the tape in a fixed radix gives polynomially many
candidate configurations, and every reachable configuration has an exact
encoding. A finite program checks an edge by comparing the local transition
and the finitely many possibly nonblank tape cells. Starting with the
initial configuration marked, repeatedly mark every successor of a marked
configuration. Each round scans polynomially many candidates; the proved
configuration bound gives a polynomial number of rounds sufficient for every
halting branch. An accepting terminal configuration is marked exactly when
the original machine accepts. A checked compiler implements the bounded
loops and table operations on a finite deterministic stack machine with
polynomial running time, giving a characteristic function in the original
definition of $\mathrm P$.
\qedhere
% lax end Lax434930Proofs.NL_subset_P
\end{proof}

\begin{lemma}\leavevmode
% lax begin Lax434930.BasicProperties
$\mathrm P\subseteq\mathrm{NP}$.
% lax end Lax434930.BasicProperties
\end{lemma}

\begin{proof}\leavevmode
% lax begin Lax434930Proofs.P_subset_NP
Let $f$ be a polynomial-time characteristic function for $A$. On an encoded input-certificate pair, a fixed finite stack machine extracts the first component in linear time. It scans the two-bit blocks before the delimiter, clears the unused suffix, and restores the extracted bits to their original order. Compose this projection with the machine for $f$. Polynomial-time composition gives a verifier language in $\mathrm P$. The zero polynomial is a certificate-length bound, and the empty certificate suffices. Decoding the first component proves the required equivalence with membership in $A$.
\qedhere
% lax end Lax434930Proofs.P_subset_NP
\end{proof}

\begin{theorem}[Savitch's theorem]\label{thm:savitch}\leavevmode
% lax begin Lax434930.PolynomialSpaceEquality
$\mathrm{PSPACE}=\mathrm{NPSPACE}$.
% lax end Lax434930.PolynomialSpaceEquality
\end{theorem}

\begin{proof}\leavevmode
% lax begin Lax434930Proofs.PSPACE_eq_NPSPACE
Every deterministic machine is also a nondeterministic machine with the same
space bound. Conversely, suppose a nondeterministic decider uses at most
$p(n)$ work cells. Enlarge the bound to the constructible function
$s(n)=p(n)+n+2$. The proved Savitch simulation gives a deterministic decider
using at most $c\,s(n)^2$ work cells, which is a polynomial bound.
\qedhere
% lax end Lax434930Proofs.PSPACE_eq_NPSPACE
\end{proof}

\begin{lemma}\leavevmode
% lax begin Lax434930.BasicProperties
$\mathrm{NP}\subseteq\mathrm{PSPACE}$.
% lax end Lax434930.BasicProperties
\end{lemma}

\begin{proof}\leavevmode
% lax begin Lax434930Proofs.NP_subset_PSPACE
Let $p$ bound the certificate length, and fix a polynomial-time verifier.
A finite nondeterministic machine constructs a unary supply of
$2p(n)+1$ symbols and guesses one bit for each symbol. Decoding the first
component of the guessed word produces a certificate of length at most
$p(n)$. Every permitted certificate occurs, since its encoding can be
padded to exactly $2p(n)+1$ bits. The machine constructs the original
input--certificate pair and simulates the verifier with finite control
and a finite work alphabet. Its stack lengths are bounded by the encoded
input length plus a constant times the verifier's running time. This is
a polynomial in $n$, and every branch halts. Thus the language belongs to
$\mathrm{NPSPACE}$. Savitch's theorem (Theorem~\ref{thm:savitch}),
proved above, gives $\mathrm{PSPACE}=\mathrm{NPSPACE}$ and hence the conclusion.
\qedhere
% lax end Lax434930Proofs.NP_subset_PSPACE
\end{proof}

\begin{lemma}\leavevmode
% lax begin Lax434930.BasicProperties
$\mathrm{NPSPACE}\subseteq\mathrm{EXPTIME}$.
% lax end Lax434930.BasicProperties
\end{lemma}

\begin{proof}\leavevmode
% lax begin Lax434930Proofs.NPSPACE_subset_EXPTIME
Apply the proved Savitch simulation to obtain a deterministic machine using
polynomial work space. A halting computation cannot revisit a configuration,
since the intervening transitions could then be repeated indefinitely.
A machine with $q$ states, $g$ work symbols, and space bound $s$ on an input
of length $n$ has at most $q(n+2)s g^s$ reachable configurations. Consequently
its running time is exponential when $s$ is polynomial. The explicit stack
simulation preserves each source transition as one stack-machine step and
includes a linear bound for clearing the stacks and writing the answer.
The stack-to-tape simulations have polynomial overhead in the input length
and simulated time, so the final single-tape bound is still exponential.
\qedhere
% lax end Lax434930Proofs.NPSPACE_subset_EXPTIME
\end{proof}

\section{An open question}
\begin{question}[P versus NP]\leavevmode
% lax begin Lax434930.PVersusNP
The P versus NP problem asks whether every language with polynomially
bounded, polynomial-time verifiable certificates can also be decided in
deterministic polynomial time. We state the conjectured separation
$\mathrm{P} \ne \mathrm{NP}$ as an open question.
% lax end Lax434930.PVersusNP
\end{question}

This open question is not used by any formal proof in the submission.
\end{document}
